% ---------------------------------------------------------------------------
% Chapter -- Nilpotent Approximation
%
% DRAFT 1 (19 Oct 2026). Proposed by Claude, to be read and rewritten in
% your own words. Every formula in this chapter was re-derived with sympy
% (code/verify_nilpotent.py) and agrees with the MATLAB code
% (one_trailer_coeff, offhooked_coeffs, x_to_z) and with
% Zuyev & Grushkovskaya, arXiv:2403.03341, Sec. III.
% The one-trailer derivation (Sec. 4.4) is your own, not in the paper.
% ---------------------------------------------------------------------------

\chapter{Nilpotent Approximation}
\label{ch:nilpotent}
\TODO[status of this chapter: about 80 percent done]

The model from Chapter~\ref{ch:model} is controllable, but its equations
contain sines and cosines. Also, no continuous feedback that depends only on the
state can stabilize it. The controllers in Chapter~\ref{ch:control} are
therefore designed for a simpler system. This system keeps only the terms that
matter most for small motions near the origin. It is called the nilpotent
approximation. This chapter explains what it is and why it is used. It is
derived for the car with one trailer (\cref{sec:approx1}) and, following
\cite{zuyev2024}, for two trailers (\cref{sec:approx2}).

% ===========================================================================
\section{Why linearization is not enough}
\label{sec:linearization}

The first idea is to linearize the system, that is, to replace it near the
origin by a linear system. The system has no drift, so the origin with zero
inputs is an equilibrium. The linearization is
\begin{equation}
  \dot x=f_1(0)\,u_1+f_2(0)\,u_2 .
  \label{eq:linearization}
\end{equation}
This system can only move in the plane spanned by $f_1(0)$ and $f_2(0)$, but the
state has five numbers (four for one trailer). So the linearization is not
controllable. It keeps the two directions that are available at once and drops
all the others. These others come from the Lie brackets, and they are exactly
the directions in which the original system is controllable
(\cref{sec:controllability}).

The fix is to look at each direction on its own scale. Suppose the inputs are
of size one and are used for a short time $\varepsilon$. Then the robot moves
about $\varepsilon$ in the first two directions, driving and turning. In a
direction that needs a manoeuvre of length $s$, it moves only about
$\varepsilon^{s}$, which is much less. So the directions cannot be measured on
the same scale. This is made exact by weights (\cref{sec:privileged}).

% ===========================================================================
\section{Privileged coordinates and weights}
\label{sec:privileged}

The construction comes from Bella\"iche, Jean and Risler~\cite{bellaiche2005},
in the form used in~\cite{zuyev2024}. It starts from the directions
$\mathcal L^s(0)$ of \cref{sec:controllability}. These are the directions at the
origin that can be made from brackets of length at most $s$, and there are $n_s$
of them. Each longer bracket adds more directions, until all of $\R^n$ is
covered. This is possible because the Lie algebra rank condition holds at the
origin. The result is a chain of spaces, each inside the next, which is called a
flag:
\begin{equation}
  \{0\}\subset\mathcal L^1(0)\subset\dots\subset\mathcal L^r(0)=\R^n ,
  \qquad \dim\mathcal L^s(0)=n_s .
\end{equation}

\paragraph{Weights.}
Every coordinate $y_j$ gets a weight. The weight is the smallest $s$ for which
some direction in $\mathcal L^s(0)$ changes $y_j$. In other words, $y_j$ has
weight $s$ if it is first reached by manoeuvres of length $s$. For the growth
vector $(2,3,4,5)$ from \eqref{eq:growth}, the weights are
\begin{equation}
  w=(1,1,2,3,4).
  \label{eq:weights}
\end{equation}
The first two coordinates have weight 1, because two directions are available at
once. The other three have weights 2, 3 and 4. For one trailer, with growth
vector $(2,3,4)$, the weights are $w=(1,1,2,3)$. In general, $w_j$ is the
smallest $s$ with $n_s\ge j$.

\paragraph{Linear step.}
Take the vectors $Y_1=f_1$, $Y_2=f_2$, $Y_3=[f_1,f_2]$, $Y_4=[f_1,[f_1,f_2]]$,
and so on. They are the columns of the matrix $F$ from \eqref{eq:Fmatrix}. At the
origin they are independent, so they can be used as new axes. The new
coordinates are
\begin{equation}
  y=F(0)^{-1}x .
  \label{eq:ylinear}
\end{equation}
The number $y_j$ says how much of the direction $Y_j(0)$ the state contains. The
vector $Y_j$ is made from brackets of length $w_j$, and so the coordinate $y_j$
has the weight $w_j$.

\paragraph{Nonlinear correction.}
The coordinates $y$ are not always good enough. The new coordinates have to be
\emph{privileged}, which means that a coordinate with weight $w_j$ changes at
most about $\varepsilon^{w_j}$ when the inputs are used for a short time
$\varepsilon$ \cite{bellaiche2005}. The coordinates $y$ do not always do this,
so they are corrected. The correction is given by the recursion
\begin{equation}
  z_q=y_q-\sum_{(\alpha,w)<w_q}\frac{1}{\alpha_1!\cdots\alpha_{q-1}!}
      \bigl(Y_1^{\alpha_1}\cdots Y_{q-1}^{\alpha_{q-1}}y_q\bigr)(0)\,
      z_1^{\alpha_1}\cdots z_{q-1}^{\alpha_{q-1}},
  \label{eq:bellaiche}
\end{equation}
where $(\alpha,w)=\sum_i w_i\alpha_i$. The recursion subtracts from $y_q$ all
terms whose weighted degree is smaller than $w_q$. For the systems in this thesis
almost all of these terms are zero. The details are in \cref{sec:approx1} and
\cref{sec:approx2}.

\paragraph{A practical test.}
Whether the coordinates are privileged can be checked on the equations. In the
exact equation for $\dot z_q$, every $z_i$ is replaced by $\varepsilon^{w_i}z_i$,
and the result is expanded in $\varepsilon$. A coordinate with weight $w_q$
changes by about $\varepsilon^{w_q}$ during a time $\varepsilon$, so its speed
$\dot z_q$ should be of size $\varepsilon^{w_q-1}$. This means that the lowest
power of $\varepsilon$ must be exactly $w_q-1$. The coefficient of this power is
the $q$-th component of the approximating system. If the lowest power is smaller,
the coordinate moves faster than its weight allows, and it needs a correction.
As an example, take $z_3=-x_2$ and $z_2=x_3$. Then
\begin{equation}
  \dot z_3=-u_1\sin x_3=-u_1\sin z_2
  \;\xrightarrow{\;z_2\to\varepsilon z_2\;}\;
  -\varepsilon\,u_1z_2+O(\varepsilon^3).
\end{equation}
The lowest power is $\varepsilon^1=\varepsilon^{w_3-1}$, because $w_3=2$. All
coordinates in this thesis were checked in this way
(\cref{app:nilpotent}).

% ===========================================================================
\section{Quasi-homogeneity}
\label{sec:homogeneity}

The weights tell how to stretch the coordinates. For a number $\lambda>0$, the
stretching is
\begin{equation}
  \delta_\lambda(z_1,\dots,z_n)=(\lambda^{w_1}z_1,\dots,\lambda^{w_n}z_n).
\end{equation}
Every coordinate is stretched by the power of $\lambda$ given by its weight. This
is the same idea as in \cref{sec:linearization}, where a short time $\varepsilon$
changes a coordinate with weight $w$ by about $\varepsilon^{w}$. The weights also
give a \emph{weighted degree} to every term. A term
$z_1^{\alpha_1}\cdots z_n^{\alpha_n}$ has the degree $\sum_iw_i\alpha_i$. For
example, $z_3$ and $z_2^2$ both have degree $2$, and $z_4$ and $z_2^3$ both have
degree $3$.

A vector field $X$ has one component $X_j$ for each coordinate. It is
\emph{weighted homogeneous of degree $m$} if every component $X_j$ has weighted
degree $w_j+m$, which means that all its terms have this degree. The fields of
the nilpotent approximation, which is derived below, have degree $-1$: the
component $j$ has degree $w_j-1$. This is what the practical test in
\cref{sec:privileged} checks. It has three consequences, and they show why the
approximation is useful.

\begin{enumerate}
  \item The bracket of two weighted homogeneous fields with degrees $m$ and $m'$
        is weighted homogeneous of degree $m+m'$. Let $g_1$ and $g_2$ be the two
        fields of the nilpotent approximation. Both have degree $-1$, so a
        bracket of length $k$ of $g_1$ and $g_2$ has degree $-k$.
  \item The fields $g_1$ and $g_2$ are polynomial, which means that their
        components are polynomials in $z$. A bracket of length $k$ has degree
        $-k$, so its component $j$ has degree $w_j-k$. A polynomial cannot have
        a negative degree, so this component is zero when $w_j<k$. Therefore all
        brackets of length greater than $r=\max_jw_j$ are zero. A Lie algebra
        with this property is called \emph{nilpotent} (of step $r$), and this is
        where the name of the approximation comes from. The exact system is not
        nilpotent, because its sines and cosines produce new terms at every
        length.
  \item A component of degree $w_j-1$ can contain only coordinates with a smaller
        weight than $w_j$. So $\dot z_j$ depends only on $z_1,\dots,z_{j-1}$, and
        the system is triangular. A triangular system can be solved one
        coordinate after the other, with a formula. The controller of
        Chapter~\ref{ch:control} uses this.
\end{enumerate}

In privileged coordinates, the exact system differs from the approximating
system only by terms of higher weighted degree~\cite{bellaiche2005}. Close to the
origin, these terms are small compared with the terms of the approximation. A
controller designed for the nilpotent approximation can therefore be expected to
work for the exact system too, as long as the difference stays small.
Chapter~\ref{ch:control} shows how this is made exact.
Chapter~\ref{ch:simulation} measures the size of the difference. There, the
controller is simulated on the approximation and on the exact kinematics, and the
two results are compared (\cref{sec:sim-exact,sec:sim-validity}).

% ===========================================================================
\section{Approximation for one trailer}
\label{sec:approx1}

In this section, the construction is carried out for the car with one off-hooked
trailer, that is, for the system with four coordinates from
\eqref{eq:onetrailer}. Its weights are $w=(1,1,2,3)$ and its degree of
nonholonomy is $r=3$. The paper~\cite{zuyev2024} treats only two trailers, so this
derivation is not in it.

\paragraph{Privileged coordinates.}
For one trailer, the matrix $F(0)$ is made of $\bar f_1$, $\bar f_2$ and their
brackets from \cref{sec:onetrailer}. It is
\begin{equation}
  F(0)=\begin{pmatrix}
    1&0&0&0\\
    0&0&-1&0\\
    0&1&0&0\\
    0&-\dfrac{\dzero}{\lone}&-\dfrac{\dzero+\lone}{\lone^{2}}&-\dfrac{\dzero+\lone}{\lone^{3}}
  \end{pmatrix}.
\end{equation}
The linear step \eqref{eq:ylinear} gives the new coordinates
\begin{equation}
  z_1=x_1,\qquad z_2=x_3,\qquad z_3=-x_2,\qquad
  z_4=\lone x_2-\frac{\dzero\lone^{2}}{\dzero+\lone}\,x_3
      -\frac{\lone^{3}}{\dzero+\lone}\,x_4 .
  \label{eq:xtoz1}
\end{equation}
They are found by writing $x$ as a mix of the four directions
$Y_1(0),\dots,Y_4(0)$ and solving for the amounts. They are called $z$, because
no nonlinear correction is needed, as shown below. The inverse map is
$x=F(0)z$, and it is used to plot the motion in the original coordinates. It
reads $x_1=z_1$, $x_2=-z_3$, $x_3=z_2$ and
$x_4=-\frac{\dzero}{\lone}z_2-\frac{\dzero+\lone}{\lone^{2}}z_3
-\frac{\dzero+\lone}{\lone^{3}}z_4$.
Only the state is written in new coordinates. The inputs $u_1$ and $u_2$ are the
physical velocities in both descriptions.

\paragraph{No nonlinear correction is needed.}
The correction \eqref{eq:bellaiche} subtracts only terms of weighted degree
smaller than $w_q$. The highest weight here is $w_4=3$, so a cubic term such as
$z_2^3$ (degree $3$) is never subtracted. In the two-trailer system it is,
because there $w_5=4>3$. Terms of degree $2$, which is less than $3$, could be
subtracted in principle, but they are not needed. The practical test below finds
no power of $\varepsilon$ below $w_q-1$. So the coordinates from the linear step
are already privileged.

\paragraph{Equations in the new coordinates.}
The first three equations follow from \eqref{eq:robot}:
$\dot z_1=u_1\cos z_2$, $\dot z_2=u_2$, $\dot z_3=-u_1\sin z_2$.
For $z_4$, the derivative of \eqref{eq:xtoz1} is computed, and $\dot x_2$,
$\dot x_3$ and $\dot x_4$ are taken from the system, using \eqref{eq:x4dot}.
This gives the exact expression
\begin{equation}
  \dot z_4=u_1\,\lone\Bigl[\sin x_3-\frac{\lone}{\dzero+\lone}\sin(x_3-x_4)\Bigr]
  -u_2\,\frac{\dzero\lone^{2}}{\dzero+\lone}\bigl[1-\cos(x_3-x_4)\bigr].
  \label{eq:z4dot}
\end{equation}
Here $x_3$ and $x_4$ have to be replaced by their expressions in $z$. From the
inverse map,
\begin{equation}
  x_3-x_4=\frac{\dzero+\lone}{\lone}\,z_2+\frac{\dzero+\lone}{\lone^{2}}\,z_3
          +\frac{\dzero+\lone}{\lone^{3}}\,z_4 .
\end{equation}
Next, every $z_i$ is replaced by $\varepsilon^{w_i}z_i$, and the result is
expanded in $\varepsilon$, as in the practical test of \cref{sec:privileged}. In
the $u_1$ term, the parts of order $\varepsilon^1$ cancel, because
$\lone\bigl[z_2-\frac{\lone}{\dzero+\lone}\cdot\frac{\dzero+\lone}{\lone}z_2\bigr]=0$.
The part of order $\varepsilon^2$ is $-z_3$. In the $u_2$ term,
$1-\cos(x_3-x_4)=\tfrac12\bigl(\tfrac{\dzero+\lone}{\lone}z_2\bigr)^2\varepsilon^2
+O(\varepsilon^3)$. Altogether
\begin{equation}
  \dot z_4=\varepsilon^2\bigl(-z_3u_1-\alpha z_2^2u_2\bigr)+O(\varepsilon^3),
  \qquad
  \alpha=\frac{\dzero(\dzero+\lone)}{2}.
  \label{eq:alpha1}
\end{equation}
The powers $\varepsilon^0$ and $\varepsilon^1$ vanish, and the lowest power is
$\varepsilon^2=\varepsilon^{w_4-1}$, as required.

\paragraph{The approximating system.}
In each equation, only the term with the lowest power of $\varepsilon$ is kept,
and $\varepsilon$ is set to $1$. This gives the nilpotent approximation for one
trailer:
\begin{equation}
  \dot z=u_1g_1(z)+u_2g_2(z),\qquad
  g_1(z)=\begin{pmatrix}1\\0\\-z_2\\-z_3\end{pmatrix},\quad
  g_2(z)=\begin{pmatrix}0\\1\\0\\-\alpha z_2^{2}\end{pmatrix}.
  \label{eq:nilp1}
\end{equation}
It has only one parameter, $\alpha$. So the system depends on the geometry only
through $\dzero$ and $\lone$, and $\alpha>0$ for all positive $\dzero$ and
$\lone$. For the platform values $\dzero=\SI{0.115}{m}$ and
$\lone=\SI{0.082}{m}$ (\cref{sec:parameters}), $\alpha\approx\SI{1.133e-2}{m^2}$.

\paragraph{Brackets and the matrix $G$.}
The brackets of the approximating fields are constant:
\begin{equation}
  [g_1,g_2]=(0,0,1,0)^\top,\qquad
  [g_1,[g_1,g_2]]=(0,0,0,1)^\top .
\end{equation}
Also, $[g_2,[g_1,g_2]]=0$, and all brackets of length four are zero. So the Lie
algebra is nilpotent of step three, as predicted in \cref{sec:homogeneity}.
The matrix $G(z)$ has the columns $g_1$, $g_2$, $[g_1,g_2]$ and
$[g_1,[g_1,g_2]]$. It is lower triangular with ones on the diagonal, so
\begin{equation}
  \det G(z)=1\qquad\text{for all }z .
  \label{eq:detG}
\end{equation}
This is different from \cref{sec:controllability}, where $\det F(0)\ne0$ was a
condition that had to be checked for the real system. Here, $\det G=1$ holds for
every system of this triangular form. It only confirms that the construction was
carried out correctly. It also means that the inverse $G(z)^{-1}$ exists for
every $z$. The controller of Chapter~\ref{ch:control} uses this inverse, so it is
well defined everywhere.

\paragraph{The chained form appears by itself.}
If $\dzero\to0$, which means that the hitch is on the axle, then $\alpha\to0$ and
\begin{equation}
  \dot z_4=-z_3u_1 .
\end{equation}
This is exactly the chained form from~\cite{sordalen1993}. The chained form was
not used anywhere in the derivation, and it comes out by itself in this case.
When the hitch is off the axle, the term $-\alpha z_2^2u_2$ is the only
difference from the classical case. This agrees with the discussion in
\cref{sec:geometry}.

% ===========================================================================
\section{Approximation for two trailers}
\label{sec:approx2}

The approximation for the car with two trailers was obtained
in~\cite{zuyev2024}. It is summarized here, together with the checks made for
this thesis.

The weights are $w=(1,1,2,3,4)$. The linear coordinates $y=F(0)^{-1}x$ are given
in~\cite{zuyev2024}, and their entries are fractions made of the four lengths.
For two trailers, the correction \eqref{eq:bellaiche} is not zero. The only term
that survives is cubic in $x_3=z_2$, and its degree is $3<w_5=4$. So
\begin{equation}
  z_1=x_1,\quad z_2=x_3,\quad z_3=-x_2,\quad z_4=y_4,\quad
  z_5=y_5-\beta x_3^{3},
  \label{eq:xtoz2}
\end{equation}
with
\begin{equation}
  \beta=\frac{\dzero\ltwo\bigl[\dzero^{2}(\done+\lone)(\lone-\ltwo-\done)
        +\lone^{3}(\dzero+\ltwo)\bigr]}{6(\dzero+\ltwo)\lone^{2}} .
\end{equation}
The coefficient $\beta$ is chosen so that it cancels the term of order
$\varepsilon^2$ in $\dot z_5$. Without it, this term would fail the practical test
of \cref{sec:privileged}.

The approximating system is
\begin{equation}
  \dot z_1=u_1,\quad \dot z_2=u_2,\quad \dot z_3=-z_2u_1,\quad
  \dot z_4=-z_3u_1-\alpha z_2^{2}u_2,\quad
  \dot z_5=(\kappa z_2^{3}-z_4)u_1-\theta z_2z_3u_2 ,
  \label{eq:nilp2}
\end{equation}
That is, $g_1=(1,0,-z_2,-z_3,\kappa z_2^3-z_4)^\top$ and
$g_2=(0,1,0,-\alpha z_2^2,-\theta z_2z_3)^\top$. The coefficients are
\begin{align}
  \alpha&=\frac{\dzero^{3}\bigl[\lone^{3}+\lone^{2}\ltwo-\lone\ltwo^{2}-\done\ltwo(\done+\ltwo)\bigr]
         +\dzero\lone^{3}\bigl[\dzero(\lone+2\ltwo)+\ltwo(\lone+\ltwo)\bigr]}
         {2(\dzero+\ltwo)\lone^{3}},\\[2pt]
  \kappa&=\frac{\dzero^{2}\Bigl\{\tfrac32\ltwo^{2}(\dzero+\lone)(\done+\lone)
         -\ltwo\bigl[\lone^{3}+\tfrac12\dzero\lone^{2}-\done\lone(\dzero+\tfrac32\done)
         -\tfrac32\dzero\done^{2}\bigr]+\tfrac12\dzero\done^{2}\lone\Bigr\}}
         {3(\dzero+\ltwo)\lone^{3}},\\[2pt]
  \theta&=\frac{\dzero^{2}\bigl[\ltwo^{2}(\done+\lone)(\dzero+\lone)
         +\bigl(\dzero\done+(\dzero+\done)\lone\bigr)\done\ltwo+\dzero\done^{2}\lone\bigr]}
         {(\dzero+\ltwo)\lone^{3}} .
\end{align}
The symbol $\alpha$ is used again, but it is not the same number as for one
trailer, because $z_4$ is defined differently here.

As a check, the two parameter sets of the numerical example in
\cite{zuyev2024} ($\dzero=\done=0.1$, $\lone=\ltwo=1$ and
$\dzero=\done=\lone=\ltwo=1$) give $(\alpha,\kappa,\theta)=(0.10495,\,0.0024,\,0.0112)$
and $(1,\,1.5,\,4)$ with these formulas. These are the values printed in the
paper. They were also checked with the MATLAB function
\texttt{offhooked\_coeffs} (Chapter~\ref{ch:feasibility}) and with a symbolic
computation (\cref{app:nilpotent}).
\TODO[the denominator of alpha in eq. (11) of the paper seems to contain $\ell_2^3$, but the derivation here gives $\ell_1^3$; the examples cannot tell because $\ell_1=\ell_2$ there; check the PDF]

The matrix $G(z)$ has the columns $g_1$, $g_2$ and the brackets up to length
four. It is again lower triangular with ones on the diagonal, so
$\det G(z)=1$. Its third column, the bracket $[g_1,g_2]$, is
$(0,0,1,0,-(\alpha+3\kappa-\theta)z_2^{2})^\top$.

Compare the approximations for one and two trailers, \eqref{eq:nilp1} and
\eqref{eq:nilp2}. The first four equations have the same form. The coefficient
$\alpha$ is different, because $z_4$ is defined differently in the two cases: for
two trailers, $y_4$ contains $x_5$. The second trailer adds one equation and a
correction $\beta$, and it gives three coefficients instead of one. It also
raises the weight of the last coordinate from $3$ to $4$.

\begin{table}[t]
  \centering
  \caption{Nilpotent approximation of the car with one and two trailers.}
  \label{tab:approx}
  \begin{tabular}{lcc}
    \toprule
                                   & one trailer & two trailers \\
    \midrule
    state dimension                & 4           & 5 \\
    growth vector                  & $(2,3,4)$   & $(2,3,4,5)$ \\
    weights                        & $(1,1,2,3)$ & $(1,1,2,3,4)$ \\
    degree of nonholonomy $r$      & 3           & 4 \\
    nonlinear correction           & none        & $\beta$ (cubic term in $z_5$) \\
    coefficients                   & $\alpha$    & $\alpha,\kappa,\theta$ \\
    $\det G(z)$                    & 1           & 1 \\
    \bottomrule
  \end{tabular}
\end{table}

% ===========================================================================
\section{Verification}
\label{sec:verification}

The derivation in \cref{sec:approx1} is not copied from a paper, so it was
checked in four ways.

\begin{enumerate}
  \item \textbf{$\det G(z)=1$.} The matrix of brackets of the approximating system
        has determinant $1$ for all $z$, for one and for two trailers. This only
        confirms that the matrix is triangular with ones on the diagonal, as the
        construction should give.
  \item \textbf{Nilpotency.} All brackets of length greater than $r$ are zero
        ($r=3$ and $r=4$), as predicted in \cref{sec:homogeneity}.
  \item \textbf{The limit $\dzero\to0$ gives the chained form.} This is a strong
        check, because the chained form was not used in the construction. It is
        a known fact from the classical theory, and the derivation reproduces it.
  \item \textbf{Symbolic computation.} The practical test of
        \cref{sec:privileged} was applied symbolically (\cref{app:nilpotent}).
        For one trailer, with $\dzero$ and $\lone$ left as symbols, it gives
        $\alpha=\dzero(\dzero+\lone)/2$. For two trailers, it determines $\beta$
        from the requirement that the $\varepsilon^2$ term of $\dot z_5$ is zero.
        It also reproduces $\alpha,\kappa,\theta,\beta$ of this chapter exactly,
        as fractions, for three sets of lengths. These are the two sets quoted in
        \cite{zuyev2024}, one of them giving $\alpha=0.10495$, $\kappa=0.0024$,
        $\theta=0.0112$, and the platform lengths.
\end{enumerate}
A general MATLAB routine, \texttt{derive\_privileged.m}, gives the same results.
It receives only the fields $f_1$ and $f_2$, and it finds the weights, the
correction and the approximating system by itself. For two trailers, it
reproduces the coefficient $\beta$ of~\cite{zuyev2024} exactly.

\paragraph{What is not verified.}
The nilpotent approximation keeps only the leading terms, so it is not exact. How
close a controller designed for it comes to the exact kinematics of
Chapter~\ref{ch:model} is measured by simulation in
Chapter~\ref{ch:simulation} (\cref{sec:sim-exact,sec:sim-validity}).
